% Options for packages loaded elsewhere
\PassOptionsToPackage{unicode}{hyperref}
\PassOptionsToPackage{hyphens}{url}
\documentclass[
  ignorenonframetext,
]{beamer}
\newif\ifbibliography
\usepackage{pgfpages}
\setbeamertemplate{caption}[numbered]
\setbeamertemplate{caption label separator}{: }
\setbeamercolor{caption name}{fg=normal text.fg}
\beamertemplatenavigationsymbolsempty
% remove section numbering
\setbeamertemplate{part page}{
  \centering
  \begin{beamercolorbox}[sep=16pt,center]{part title}
    \usebeamerfont{part title}\insertpart\par
  \end{beamercolorbox}
}
\setbeamertemplate{section page}{
  \centering
  \begin{beamercolorbox}[sep=12pt,center]{section title}
    \usebeamerfont{section title}\insertsection\par
  \end{beamercolorbox}
}
\setbeamertemplate{subsection page}{
  \centering
  \begin{beamercolorbox}[sep=8pt,center]{subsection title}
    \usebeamerfont{subsection title}\insertsubsection\par
  \end{beamercolorbox}
}
% Prevent slide breaks in the middle of a paragraph
\widowpenalties 1 10000
\raggedbottom
\AtBeginPart{
  \frame{\partpage}
}
\AtBeginSection{
  \ifbibliography
  \else
    \frame{\sectionpage}
  \fi
}
\AtBeginSubsection{
  \frame{\subsectionpage}
}
\usepackage{iftex}
\ifPDFTeX
  \usepackage[T1]{fontenc}
  \usepackage[utf8]{inputenc}
  \usepackage{textcomp} % provide euro and other symbols
\else % if luatex or xetex
  \usepackage{unicode-math} % this also loads fontspec
  \defaultfontfeatures{Scale=MatchLowercase}
  \defaultfontfeatures[\rmfamily]{Ligatures=TeX,Scale=1}
\fi
\usepackage{lmodern}
\ifPDFTeX\else
  % xetex/luatex font selection
\fi
% Use upquote if available, for straight quotes in verbatim environments
\IfFileExists{upquote.sty}{\usepackage{upquote}}{}
\IfFileExists{microtype.sty}{% use microtype if available
  \usepackage[]{microtype}
  \UseMicrotypeSet[protrusion]{basicmath} % disable protrusion for tt fonts
}{}
\makeatletter
\@ifundefined{KOMAClassName}{% if non-KOMA class
  \IfFileExists{parskip.sty}{%
    \usepackage{parskip}
  }{% else
    \setlength{\parindent}{0pt}
    \setlength{\parskip}{6pt plus 2pt minus 1pt}}
}{% if KOMA class
  \KOMAoptions{parskip=half}}
\makeatother
\usepackage{longtable,booktabs,array}
\usepackage{caption}
\captionsetup[table]{skip=6pt}
\newcounter{none} % for unnumbered tables
\usepackage{calc} % for calculating minipage widths
\usepackage{caption}
% Make caption package work with longtable
\makeatletter
\def\fnum@table{\tablename~\thetable}
\makeatother
\setlength{\emergencystretch}{3em} % prevent overfull lines
\providecommand{\tightlist}{%
  \setlength{\itemsep}{0pt}\setlength{\parskip}{0pt}}
% Auto-generated by infrastructure/rendering/_pdf_latex_helpers.py::extract_math_font_preamble.
% Minimal header passed to Pandoc via -H so Beamer slides pick up the same math font
% as the combined-PDF path. Adjust the manuscript's preamble.md to change the font globally.
\usepackage{unicode-math}
\setmathfont{latinmodern-math.otf}

% Manuscript macro/package fallbacks (newcommand -> providecommand).
% Core mathematics
\usepackage{amsmath}
\usepackage{amssymb}
% Document layout
% Tables
\usepackage{booktabs}
% Algorithm typesetting (for pseudocode in §2 Methodology)
% Code listings
\usepackage{listings}
% Typography and formatting
% Cross-references and citations
% ── Unicode-capable mono font for code listings ──────────────────────
% JuliaMono (TeX Live 2026) covers the full math/Greek glyph set used in
% scientific code blocks; the default lmmono lacks \alpha/\mu/\partial/\nabla.
%
% The hardcoded TeX Live 2026 path below is portable to most macOS / Linux
% TeX Live installs that placed JuliaMono under the default texmf-dist path.
% If your TeX Live version differs (e.g., 2025, 2027) or your distribution
% installs fonts elsewhere, comment out the explicit Path = line — fontspec
% will then search the system font path. If JuliaMono is not installed at
% all, comment out the whole block; xelatex falls back to lmmono with
% reduced Unicode coverage but the document still compiles.
% Math font for unicode-math: Latin Modern Math (TeX Live) has full BMP
% coverage including U+2223 (\mid), U+226A/226B (\ll/\gg), and the Greek/
% blackboard letters used in equations. Without an explicit \setmathfont,
% unicode-math falls back to lmroman text font which lacks several glyphs
% and emits "Missing character" warnings on every \mid in math mode.

% Slide layout defaults for warning-clean scientific decks.
\usepackage{etoolbox}
\IfFileExists{xurl.sty}{\usepackage{xurl}}{}
\IfFileExists{seqsplit.sty}{\usepackage{seqsplit}}{\newcommand{\seqsplit}[1]{#1}}
\protected\def\breakseq#1{\seqsplit{#1}}
\protected\def\breaktt#1{\begingroup\ttfamily\seqsplit{#1}\endgroup}
\setlength{\emergencystretch}{6em}
\tolerance=5000
\hbadness=10000
\hfuzz=1pt
\setlength{\tabcolsep}{2pt}
\AtBeginEnvironment{longtable}{\tiny\renewcommand{\arraystretch}{0.86}\setlength{\tabcolsep}{1pt}}
\AtBeginEnvironment{tabular}{\tiny\renewcommand{\arraystretch}{0.86}\setlength{\tabcolsep}{1pt}}
\AtBeginEnvironment{equation}{\tiny}
\AtBeginEnvironment{equation*}{\tiny}
\AtBeginEnvironment{align}{\tiny}
\AtBeginEnvironment{align*}{\tiny}
\AtBeginEnvironment{itemize}{\footnotesize}
\AtBeginEnvironment{enumerate}{\footnotesize}
\AtBeginEnvironment{description}{\footnotesize}
\setbeamerfont{caption}{size=\tiny}
\setbeamerfont{caption name}{size=\tiny}
\setbeamerfont{normal text}{size=\small}
\setbeamerfont{frametitle}{size=\small}
\setbeamerfont{section title}{size=\footnotesize}
\setbeamerfont{subsection title}{size=\footnotesize}
\setbeamertemplate{section page}{%
  \centering
  \begin{beamercolorbox}[sep=12pt,center,wd=\paperwidth]{section title}
    \parbox{0.86\paperwidth}{\centering\usebeamerfont{section title}\insertsection\par}
  \end{beamercolorbox}
}
\setbeamertemplate{subsection page}{%
  \centering
  \begin{beamercolorbox}[sep=8pt,center,wd=\paperwidth]{subsection title}
    \parbox{0.86\paperwidth}{\centering\usebeamerfont{subsection title}\insertsubsection\par}
  \end{beamercolorbox}
}
\setlength{\abovecaptionskip}{2pt}
\setlength{\belowcaptionskip}{0pt}

% Natbib and cross-reference fallbacks — slides don't load natbib
% or cleveref, but combined-PDF manuscript prose may emit these
% commands. The fallback renders citations as a bracketed key list
% and cross-references as detokenized labels so slides don't fail on
% undefined control sequences or raw underscores. \providecommand is
% a no-op if packages load later.
\providecommand{\citep}[1]{[#1]}
\providecommand{\citet}[1]{#1}
\providecommand{\citealp}[1]{#1}
\providecommand{\citeauthor}[1]{#1}
\providecommand{\citeyear}[1]{#1}
\providecommand{\cref}[1]{\texttt{\detokenize{#1}}}
\providecommand{\Cref}[1]{\texttt{\detokenize{#1}}}
\providecommand{\crefrange}[2]{\texttt{\detokenize{#1}}--\texttt{\detokenize{#2}}}
\providecommand{\Crefrange}[2]{\texttt{\detokenize{#1}}--\texttt{\detokenize{#2}}}
\providecommand{\paragraph}[1]{\textbf{#1}\ }

\newtheorem{proposition}{Proposition}
\newtheorem{hypothesis}{Hypothesis}
\newtheorem{remark}[theorem]{Remark}
\newtheorem{axiom}{Axiom}
\newtheorem{property}{Property}
\usepackage{bookmark}
\IfFileExists{xurl.sty}{\usepackage{xurl}}{} % add URL line breaks if available
\urlstyle{same}
% fallback for those not using the hyperref driver hyperxmp:
\makeatletter
\@ifundefined{xmpquote}{\newcommand{\xmpquote}[1]{#1}}{}
\makeatother
\hypersetup{
  hidelinks,
  pdfcreator={LaTeX via pandoc}}

\author{\texorpdfstring{}{}}
\date{}

\begin{document}

\section{Scenario Recommendations with Change Conditions and Confidence
Tiers}\label{sec:scenarios}

\begin{frame}[allowframebreaks]{Scenario Recommendations with Change
Conditions and Confidence Tiers}
The 8 scenarios below map the evaluation of
sec.~4 onto recurring situations.
Recommendations are conditional judgments from documented designs and
advisories, not an absolute ordering across every threat model. The
assumptions matter enough that changing the hardware, the required
software, or the adversary can change the recommendation; the change
conditions state what would.
\end{frame}

\subsection{Scenario recommendations}\label{scenario-recommendations}

\begin{frame}[allowframebreaks]{Scenario recommendations}
tbl.~\ref{tbl:scenarios} compresses the recommendations. Each
recommended direction names the architectural property doing the work,
not merely the distribution label, because the properties ---
containment, authority, integrity, update operations, and the rest of
the 9 (tbl.~2) --- are what survive a change of
vendor.

\end{frame}

\begin{frame}[allowframebreaks]{Scenario recommendations}
\begin{longtable}[]{@{}
  >{\raggedright\arraybackslash}p{(\linewidth - 4\tabcolsep) * \real{0.3333}}
  >{\raggedright\arraybackslash}p{(\linewidth - 4\tabcolsep) * \real{0.3333}}
  >{\raggedright\arraybackslash}p{(\linewidth - 4\tabcolsep) * \real{0.3333}}@{}}
\caption{The 8 scenario recommendations: the architectural property
doing the work, the direction, and the condition that would change
it}\label{tbl:scenarios}\tabularnewline
\toprule\noalign{}
\begin{minipage}[b]{\linewidth}\raggedright
Scenario
\end{minipage} & \begin{minipage}[b]{\linewidth}\raggedright
Recommended direction
\end{minipage} & \begin{minipage}[b]{\linewidth}\raggedright
Condition that could change it
\end{minipage} \\
\midrule\noalign{}
\endfirsthead
\toprule\noalign{}
\begin{minipage}[b]{\linewidth}\raggedright
Scenario
\end{minipage} & \begin{minipage}[b]{\linewidth}\raggedright
Recommended direction
\end{minipage} & \begin{minipage}[b]{\linewidth}\raggedright
Condition that could change it
\end{minipage} \\
\midrule\noalign{}
\endhead
High-risk personal workstation with distinct sensitive identities &
Qubes, with deliberately separated domains, disposable intake, minimal
administration, and prompt updates; release 4.3.0's continued
disaggregation of the GUI and administrative stacks makes the
compartmentalized workflow easier to keep (sec.~5) &
Incompatible hardware, essential unsupported software, or inability to
sustain the compartmentalized workflow \\
Conventional Linux desktop with a hardening priority & Evaluate
secureblue first against actual application requirements --- version
4.9.1 ships a SUID-less baseline with SELinux-enforced namespace
restrictions; Fedora Workstation or an atomic Fedora desktop when a more
standard environment is more sustainable (sec.~7) &
Hardening restrictions require broad exceptions, or project maintenance
does not meet the owner's assurance requirements \\
Skilled operator prioritizing auditable configuration and replaceable
environments & NixOS, paired with explicit runtime confinement,
boot-integrity decisions, runtime secret handling, and a disciplined
update process; the project's hardening wiki now documents a coherent
baseline after the hardened profile's removal, and the provenance
tooling of sec.~14 strengthens the
audit story (sec.~6) & The custom security integration
becomes too complex to review and maintain \\
Autonomous coding or research agents processing hostile material &
Disposable task VMs or comparable isolated workers, with controlled
egress and externally mediated credentials; Nix can help construct the
workers (sec.~8) & Tasks demand high-impact access that
cannot be safely scoped; those actions remain separately authorized \\
Controlled Kubernetes node fleet & Talos; also assess Bottlerocket where
its container-host ecosystem fits; Talos 1.12's configurable TPM PCR
policies make its Secure Boot posture declarable per fleet
(sec.~8) & Orchestrator, hardware, extension, or
operational requirements conflict with the reduced host interface \\
General enterprise server & RHEL or Ubuntu LTS with explicit support
scope and tested policies; Debian or NixOS where the operator can own
the maintenance obligations (sec.~8) & Package coverage,
response requirements, or configuration expertise favor a different
support model \\
Anonymity or reduced local traces & Whonix for persistent separated
work; Tails for amnesic sessions (sec.~9) & The
actual priority is protecting local credentials from an
already-compromised application rather than anonymity \\
High-assurance architectural research & Track seL4 and Genode/Sculpt
while evaluating the complete deployed system and its proof boundary
(sec.~9) & Production compatibility, drivers,
tooling, or operational assurance outweighs theoretical elegance \\
\bottomrule\noalign{}
\end{longtable}

\end{frame}

\begin{frame}[allowframebreaks]{Scenario recommendations}
Two rows deserve emphasis. Autonomous agent hosting is the scenario the
threat model of sec.~3 is built around: the
recommendation is disposable isolation plus externally mediated
credentials, not a host choice, and no reviewed distribution turns an
ordinary container or VM into that boundary by naming. Anonymity is
deliberately narrow: Whonix's gateway/workstation model and Tails'
amnesic design address trace reduction and session separation (W.
Project 2026a; T. Project 2026), and neither makes an authorized agent
harmless --- protecting local credentials from an already-compromised
application is the workstation containment problem, not the anonymity
problem (W. Project 2026b). Matching scenario to priority is the whole
recommendation.

\end{frame}

\begin{frame}[allowframebreaks]{Scenario recommendations}
The refresh behind these rows is current, and it is worth stating what
changed without overstating it. Qubes 4.3.0 reached general availability
with the GUI stack and administrative services disaggregated out of
dom0, continuing the small-privileged-interface trajectory (Project
2025). secureblue's v4.9.1 ships a SUID-less userland with
SELinux-enforced namespace restrictions, extending its hardening line
from allocator and browser work into removing an entire privilege class
(secureblue Project 2026b). NixOS 26.05 ships without the long-debated
hardened profile --- removed from nixpkgs rather than merely deprecated
({community} 2026b) --- with the NixOS Hardening wiki now serving as the
documented baseline an operator assembles from ({community} 2026a).
Talos 1.12 made its TPM measurement policies configurable, so a fleet
can state which PCRs seal what (Labs 2025). Each release lowers the
\end{frame}

\begin{frame}[allowframebreaks]{Scenario recommendations}
operational cost of keeping a recommendation that was already the
architectural judgment; none of them changes a direction, and none
converts a moderate row into a high one.
\end{frame}

\subsection{The strongest overall
direction}\label{the-strongest-overall-direction}

\begin{frame}[allowframebreaks]{The strongest overall direction}
The strongest synthesis across the reviews is: \textbf{Qubes-like
containment, Nix-like reproducible operations, strong boot and update
integrity, and capability-limited agents.} This is a design target, not
a product claim. No reviewed system ships all four properties as a
default, and they are exactly the non-substitutable set that
sec.~4 argues must be evaluated separately:
containment without reproducible operations decays; reproducible
operations without containment standardizes a single wide domain; both
without boot and update integrity rest on an unverified foundation; all
three without capability-limited agents leave the authorized-misuse path
open (sec.~10). The trust-domain architecture
of tbl.~8 is the operating plan for the target,
and the operator and orchestration disciplines of sec.~12
and sec.~13 are what keep it standing during real
work.

\end{frame}

\begin{frame}[allowframebreaks]{The strongest overall direction}
Operator readiness is trainable as well as structural:
analytic-tradecraft curricula for agentic work --- the educational line
behind the inspectable skill libraries cited in
sec.~11 --- turn the operator's own review
judgment into a maintained control, on the same footing as the
configurations and rehearsed responses the recommendations above assume
(Friedman 2026).

The design-target framing carries a practical corollary: a
well-maintained implementation of part of this design can be safer than
an ambitious combination whose integration nobody reliably owns. The
Qubes-plus-NixOS combination is the standing example. Community work on
a NixOS template exists (Q. O. Project 2026b), but community templates
do not receive updates from the Qubes project itself (Q. O. Project
2026c), and the integration --- not either component --- is where
maintenance risk accumulates (sec.~6). Ownership is part
\end{frame}

\begin{frame}[allowframebreaks]{The strongest overall direction}
of the security property: who reviews the composition, who patches it,
who rehearses its recovery. An unowned maximal stack decays into exactly
the neglected-architecture failure that sec.~5 and
sec.~6 both document, while a smaller, owned one compounds
its advantages over time.
\end{frame}

\subsection{Confidence tiers and evidence
limits}\label{confidence-tiers-and-evidence-limits}

\begin{frame}[allowframebreaks]{Confidence tiers and evidence limits}
tbl.~\ref{tbl:confidence} states what this review claims, with what
confidence, and what limits the evidence behind each claim.

\end{frame}

\begin{frame}[allowframebreaks]{Confidence tiers and evidence limits}
\begin{longtable}[]{@{}
  >{\raggedright\arraybackslash}p{(\linewidth - 4\tabcolsep) * \real{0.3333}}
  >{\raggedright\arraybackslash}p{(\linewidth - 4\tabcolsep) * \real{0.3333}}
  >{\raggedright\arraybackslash}p{(\linewidth - 4\tabcolsep) * \real{0.3333}}@{}}
\caption{The confidence tiers and evidence limits attached to this
review's principal claims}\label{tbl:confidence}\tabularnewline
\toprule\noalign{}
\begin{minipage}[b]{\linewidth}\raggedright
Confidence tier
\end{minipage} & \begin{minipage}[b]{\linewidth}\raggedright
Claim
\end{minipage} & \begin{minipage}[b]{\linewidth}\raggedright
Evidence basis and limits
\end{minipage} \\
\midrule\noalign{}
\endfirsthead
\toprule\noalign{}
\begin{minipage}[b]{\linewidth}\raggedright
Confidence tier
\end{minipage} & \begin{minipage}[b]{\linewidth}\raggedright
Claim
\end{minipage} & \begin{minipage}[b]{\linewidth}\raggedright
Evidence basis and limits
\end{minipage} \\
\midrule\noalign{}
\endhead
High & The properties are independent: build reproducibility, boot
integrity, runtime isolation, anonymity, and agent authorization answer
different questions, and no sound selection substitutes one for another
& Follows from the property definitions of
sec.~4; structural, not an empirical
measurement \\
High & Excessively broad legitimate access is dangerous independently of
exploit resistance; the AISI incident is the key precedent --- its own
report distinguishes permitted internet activity from sandbox escape,
with out-of-scope actions in 10 of 122 runs inside an intact sandbox and
no resulting real-world harm found (Institute 2026) & A primary incident
report; its authors note the configurations were not representative of
public access --- the lesson is structural, not a base rate \\
Moderate & Qubes is the strongest fit among the reviewed deployable
workstation candidates for this compartmentalization-centered threat
model, under substantial operational conditions (Q. O. Project 2026a) &
An architectural assessment from documented design and advisories, not a
measured probability of resisting any particular attacker \\
Moderate & secureblue is a particularly interesting conventional-desktop
candidate, and NixOS a particularly strong compositional platform,
pending independent audit coverage and patch-latency evidence
(secureblue Project 2026a; N. Project 2026) & Project-documented
features; public descriptions do not settle comparative implementation
quality, audit coverage, or time-to-patch \\
Low & Any unconditional prediction of which named distribution will be
most secure across 2028--2031 & Declined: project execution, hardware
changes, ecosystem support, and unknown vulnerabilities can reorder the
field; the review does not manufacture a winner \\
\bottomrule\noalign{}
\end{longtable}

\end{frame}

\begin{frame}[allowframebreaks]{Confidence tiers and evidence limits}
The high-confidence rows are structural. Property independence follows
from the definitions, and the authorized-misuse danger is demonstrated
rather than projected: the AISI report is an incident report by the
organization that ran the testing, and its value is precisely that it
documents unsanctioned agent behavior without a sandbox escape
(Institute 2026). The moderate rows are architectural judgments with
conditions attached, not probabilities; they weaken under the stated
change conditions rather than under argument. It is worth noting where
official posture now sits relative to that moderation: the CISA-led Five
Eyes guidance on careful adoption of agentic AI services recommends
sandboxed deployment, staged low-risk-task adoption, threat-model-based
evaluation, and red-teaming (Cybersecurity and Agency 2026) --- a
guarded-adoption posture convergent with, though independent of, this
\end{frame}

\begin{frame}[allowframebreaks]{Confidence tiers and evidence limits}
review's moderate tiers. Convergence is not corroboration of the
specific platform judgments, which remain assessments from documented
design; it does indicate that the posture itself is no longer
idiosyncratic. The low-confidence row is a refusal, consistent with the
forecast posture of sec.~15.
\end{frame}

\begin{frame}[allowframebreaks]{Evidence limits}
\protect\phantomsection\label{evidence-limits}
The review relied principally on official documentation, project
security advisories, and primary incident or research reports, with
developer discussions used for specific integration caveats. It did not
install the candidates, formally audit their policies, verify the
reproducibility of their images, or run offensive-agent trials against
them. Where current defaults or security properties were not verified,
the comparison labels them n.a. or withholds the stronger claim --- as
with the unverified role-specific defaults noted for MicroOS in
sec.~8 and the unverified boot-integrity baseline noted
for Kicksecure in sec.~7. Absence of a confirmed
feature in this review is not evidence that the feature is unavailable;
it is a statement about this review's evidence. Readers extending this
work should treat the scenario table as a starting position to be
\end{frame}

\begin{frame}[allowframebreaks]{Evidence limits}
re-derived under their own assumptions, using the properties of
sec.~4 rather than distribution labels.
\end{frame}

\end{document}
